\documentclass[11pt,a4paper]{article}

\usepackage{fontspec}
\setmonofont{JetBrainsMono Nerd Font Mono}[Scale=0.78]
\usepackage{amsmath,amssymb,amsthm}
\usepackage{geometry}
\geometry{margin=2.5cm}
\usepackage{booktabs}
\usepackage{listings}
\usepackage{xcolor}
\usepackage{hyperref}

\hypersetup{
  colorlinks=true,
  linkcolor=blue,
  citecolor=blue,
  urlcolor=blue,
  pdftitle={CollatzAdicBase: Machine-Verified 3-Adic Rational Descent and 2-Adic Glide Bounds for the Collatz Map in Lean 4},
  pdfauthor={Jovian Wilson Simon}
}

\lstdefinestyle{lean}{
  basicstyle=\ttfamily\footnotesize,
  columns=fullflexible,
  keepspaces=true,
  breaklines=true,
  frame=single,
  framesep=6pt,
  xleftmargin=6pt,
  xrightmargin=6pt,
  backgroundcolor=\color{gray!6}
}
\lstset{style=lean}

\newtheorem{theorem}{Theorem}[section]
\newtheorem{lemma}[theorem]{Lemma}
\newtheorem{proposition}[theorem]{Proposition}
\newtheorem{corollary}[theorem]{Corollary}
\newtheorem{definition}[theorem]{Definition}

\title{\textbf{CollatzAdicBase: Machine-Verified 3-Adic Rational Descent and 2-Adic Glide Bounds for the Collatz Map in Lean 4}}

\author{Jovian Wilson Simon\\
A J Institute of Engineering and Technology\\
ORCID: 0009-0000-5022-8570}

\date{November 2026}

\begin{document}
\maketitle

\begin{abstract}
We report a complete machine-checked formalization, in Lean 4 with Mathlib
(version 4.36.0), of the local $2$-adic and $3$-adic valuation dynamics that
underlie two earlier preprints on the Collatz map.  On the $3$-adic side we
formalize reduced rational states $(a,d)$, prove that the denominator gcd at an
odd step divides $3$, establish stability of every prime $p\neq 3$, prove that
an odd step removes exactly one factor of $3$ when $3\mid d$, and derive a
descent theorem for pure powers of three together with a kernel-checked
exhaustive verification for $1\le k\le 6$.  On the $2$-adic side we formalize
the glide predicates and prove both the initial fuel law
$G(2^{k}q-1)\ge 2k+1$ and the alternating-bit quantization
$G\bigl((2^{2k+3}+1)/3\bigr)=8$ for all $k\ge 1$.  The formalization also
uncovers an off-by-one error in Lemma 3 of the second preprint, which we correct
and document.  A full \texttt{\#print axioms} audit of all five verified modules
shows that every theorem depends only on the kernel's core axioms
(\texttt{propext}, \texttt{Classical.choice}, \texttt{Quot.sound}); there is no
\texttt{sorryAx} and no project-specific axiom.  The repository builds with zero
errors and zero warnings.  This paper is the third in the series and is intended
to accompany the source release \texttt{v1.0.0}.
\end{abstract}

\section{Introduction}
\label{sec:intro}

The Collatz map
\[
  T(n)=\begin{cases}
    n/2, & n\equiv 0 \pmod 2,\\
    3n+1, & n\equiv 1 \pmod 2,
  \end{cases}
\]
and its extension to rationals with odd denominator occupy a peculiar position
in arithmetic dynamics: the global conjecture remains open, yet many of the
\emph{local} mechanisms that drive individual orbits are elementary and
explicit.  The two preceding preprints in this series studied those mechanisms
from complementary directions.  The first~\cite{paper1} analyzed periodic
orbits of the $2$-adic rational Collatz map, proving $3$-adic exclusivity of
nontrivial cycles and a descent theorem on pure $3$-power denominators.  The
second~\cite{paper2} studied the integer map up to $10^{7}$, proving a rigid
glide window for the alternating-bit family and a general $2$-adic ``fuel law''
for odd integers with many trailing one-bits.

The purpose of the present paper is to place the \emph{local} core of both
preprints on an uncompromising foundation: a proof assistant kernel.  We do not
claim a proof of the Collatz conjecture.  We claim something narrower but fully
verified: every valuation-theoretic step used in~\cite{paper1,paper2} is a
theorem of Lean 4, checked from the axioms of the kernel alone.

Mechanization matters here for three reasons.  First, the arguments are replete
with side conditions (positivity, coprimality, parity, subtraction truncation in
$\mathbb{N}$) that are easy to suppress in informal prose but fatal to a proof;
the kernel forces them into the open.  Second, the exhaustive $k\le 6$
computation is a genuine finite verification, not a heuristic: it is performed
by reduction of a decidable Boolean and accepted by the same kernel.  Third,
as we describe in Section~\ref{sec:corrigendum}, the mechanization detected a
concrete off-by-one error in the published statement of Lemma 3 of~\cite{paper2}.

The paper is organized as follows.  Section~\ref{sec:arch} documents the
formalization environment and the state representations.  Section~\ref{sec:descent}
presents the $3$-adic rational descent.  Section~\ref{sec:glide} presents the
$2$-adic fuel law and alternating glide bounds.  Section~\ref{sec:corrigendum}
gives the verified corrigendum.  Section~\ref{sec:audit} reports the axiom audit
and exact reproduction instructions.

\section{Formalization architecture and environment}
\label{sec:arch}

\subsection{Toolchain and project structure}

The development is a Lean 4 project pinned to toolchain and Mathlib version
\texttt{v4.36.0}.  It comprises one root module \texttt{CollatzAdicBase} and five
verified submodules, all exported together:

\begin{center}
\begin{tabular}{@{}ll@{}}
\toprule
Module & Content \\
\midrule
\texttt{Basic} & Collatz/shortcut operators, iterate lemmas, glide predicates \\
\texttt{FuelLaw} & Proposition 4: initial $2$-adic fuel law \\
\texttt{RationalDescent} & $3$-adic rational descent, Lemmas 2.1--2.4, Theorem 1, small-$k$ check \\
\texttt{Refueling} & Mid-flight refueling congruence (Paper 2, Section 6) \\
\texttt{AlternatingGlide} & Theorem 1, alternating-bit glide, Lemma 3 \\
\texttt{LatentSaturation} & Denominator lifting and cycle saturation (Paper 4) \\
\texttt{Axioms} & \texttt{\#print axioms} audit over all five verified modules \\
\bottomrule
\end{tabular}
\end{center}

\noindent The root file \texttt{CollatzAdicBase.lean} imports all seven modules,
and \texttt{CollatzAdicBase/Axioms.lean} re-audits the five
theorem-bearing ones.
The library depends only on Lean's core library and Mathlib; it declares no
custom axioms and uses no \texttt{sorry}.

\subsection{State representations and operators}

A reduced positive rational with odd denominator is encoded as a pair of
naturals \texttt{(a,d)} with $a>0$, $d>0$, $d$ odd, and $\gcd(a,d)=1$.  In Lean:

\begin{lstlisting}
abbrev State := ℕ × ℕ

def Valid (x : State) : Prop :=
  0 < x.1 ∧ 0 < x.2 ∧ Odd x.2 ∧ Nat.Coprime x.1 x.2

def step (x : State) : State :=
  if _ : Even x.1 then (x.1 / 2, x.2)
  else
    let g := Nat.gcd (3 * x.1 + x.2) x.2
    ((3 * x.1 + x.2) / g, x.2 / g)
\end{lstlisting}

\noindent Thus an even numerator triggers a halving that leaves the denominator
unchanged; an odd numerator triggers the affine step $a\mapsto 3a+d$ followed by
cancellation of the exact gcd $g=\gcd(3a+d,d)$.  Iteration is written with
Mathlib's \texttt{Nat.iterate}: \texttt{orbit x n := step\textasciicircum[n] x}.  The integer
operators are defined in \texttt{Basic.lean}:

\begin{lstlisting}
def collatz (n : Nat) : Nat := if n % 2 = 0 then n / 2 else 3 * n + 1

def shortcut (m : Nat) : Nat := (3 * m + 1) / 2

def GlideAtLeast (n m : Nat) : Prop :=
  ∀ t, 1 ≤ t → t ≤ m → Nat.iterate collatz t n > n

def GlideEq (n m : Nat) : Prop :=
  (∀ t, 1 ≤ t → t < m → Nat.iterate collatz t n > n) ∧
  Nat.iterate collatz m n < n
\end{lstlisting}

\noindent Here \texttt{shortcut} compresses one odd step and its forced even
successor; \texttt{GlideAtLeast n m} says the orbit stays strictly above $n$
through step $m$, and \texttt{GlideEq n m} says it stays strictly above through
step $m-1$ and drops strictly below at step $m$.  In the notation of~\cite{paper2},
$G(n)\ge m+1$ in the first case and $G(n)=m$ in the second.

\subsection{Kernel-checked computation}

Finite verification uses the reflexivity of equality on closed Boolean terms.
A proposition $P$ is certified by producing a term of $P$ obtained from a
\texttt{decide} proof of a Boolean reduction; the kernel then checks the
reduction.  For the small-$k$ computation we use two layers: a Boolean sieve
\texttt{verifyPureThreePowerSmallK} proved equal to \texttt{true} by
\texttt{decide}, and a propositional corollary
\texttt{verification\_all\_reduced\_small\_k} obtained by reflecting the sieve
through \texttt{reachesOneBool\_iff}.  No external oracle or trusted computation
is involved.

\section{Mechanized 3-adic rational descent}
\label{sec:descent}

We present the Lean signatures verbatim, followed by proof outlines.  All
statements in this section live in namespace \texttt{CollatzAdicBase} and come
from \texttt{CollatzAdicBase/RationalDescent.lean}.

\subsection{The gcd lemma}

\begin{lemma}[gcd divides three]
\label{thm:gcd3}
\begin{lstlisting}
lemma gcd_dvd_three {a d : ℕ} (hcop : Nat.Coprime a d) :
    Nat.gcd (3 * a + d) d ∣ 3
\end{lstlisting}
\end{lemma}

\begin{proof}[Proof outline]
Put $g=\gcd(3a+d,d)$.  Then $g\mid d$ and $g\mid 3a+d$, hence by reduction modulo
$g$ we get $g\mid 3a$.  Since $g\mid d$ and $\gcd(a,d)=1$, we have $\gcd(g,a)=1$.
Euclid's lemma applied to $g\mid 3a$ then gives $g\mid 3$.
\end{proof}

Since $3$ is prime, \texttt{dvd\_three\_cases} converts $g\mid 3$ into
$g=1\vee g=3$; this dichotomy drives the rest of the section.

\subsection{Stability of primes other than three}

\begin{lemma}[Lemma 2.1]
\label{thm:stab}
\begin{lstlisting}
lemma lemma_2_1_stability (p : ℕ) (hp : p.Prime) (hp3 : p ≠ 3)
    {a d : ℕ} (hcop : Nat.Coprime a d) :
    padicValNat p (step (a, d)).2 = padicValNat p d
\end{lstlisting}
\end{lemma}

\begin{proof}[Proof outline]
Even steps leave $d$ untouched.  For odd steps, write $g=\gcd(3a+d,d)$ and use
Lemma~\ref{thm:gcd3}: either $g=1$, which is trivial, or $g=3$, in which case
$g$ has $p$-adic valuation $0$ because $p\neq 3$.  Rewriting
$\operatorname{padicValNat}$ as a factorization and applying
\texttt{Nat.factorization\_div} gives the claim.
\end{proof}

\subsection{One factor of three per odd step}

\begin{lemma}[Lemma 2.2]
\label{thm:one3}
\begin{lstlisting}
lemma lemma_2_2_one_three_per_odd_step {a d : ℕ} (_ha : Odd a) (_hdpos : 0 < d)
    (hcop : Nat.Coprime a d) (h3d : 3 ∣ d) :
    Nat.gcd (3 * a + d) d = 3 ∧
      padicValNat 3 (d / 3) = padicValNat 3 d - 1
\end{lstlisting}
\end{lemma}

\begin{proof}[Proof outline]
With $g=\gcd(3a+d,d)$, Lemma~\ref{thm:gcd3} gives $g\mid 3$, while $3\mid d$
and $3\mid 3a$ give $3\mid g$; hence $g=3$.  The valuation identity follows from
\texttt{Nat.factorization\_div} and the fact that the factorization of $3$ at
the prime $3$ equals $1$.
\end{proof}

\subsection{Forced halving after an odd step}

\begin{lemma}[Lemma 2.3]
\label{thm:halving}
\begin{lstlisting}
lemma lemma_2_3_forced_halving {a d : ℕ} (ha : Odd a) (hd : Odd d)
    (hdpos : 0 < d) (_hcop : Nat.Coprime a d) :
    Even ((3 * a + d) / Nat.gcd (3 * a + d) d)
\end{lstlisting}
\end{lemma}

\begin{proof}[Proof outline]
Let $g=\gcd(3a+d,d)$.  Because $d$ is odd, $g$ is odd.  The numerator $3a+d$ is
even (odd $+$ odd), so $3a+d=2m$; since $g\mid 2m$ and $g$ is odd, $g\mid m$.
Writing $m=gt$ gives $3a+d=g(2t)$, hence $(3a+d)/g=2t$ by
\texttt{Nat.mul\_div\_right}.
\end{proof}

\subsection{Cycles are three-free}

\begin{lemma}[Lemma 2.4 and Corollary 2.5]
\label{thm:cycles}
\begin{lstlisting}
lemma lemma_2_4_and_cor_2_5_three_free_cycles {x : State} (hx : Valid x) {L : ℕ}
    (hL : 1 ≤ L) (hcycle : orbit x L = x) :
    ¬ 3 ∣ x.2 ∧ ∀ i, i ≤ L → (orbit x i).2 = x.2
\end{lstlisting}
\end{lemma}

\begin{proof}[Proof outline]
First, an all-even cycle is impossible: \texttt{orbit\_allEven\_decreases} shows
that $L$ consecutive even steps strictly decrease the numerator, contradicting
\texttt{orbit x L = x}.  Hence there is a first odd-numerator index $i_{0}$; by
minimality the denominator is still $d$ at $i_{0}$.  If $3\mid d$, Lemma~\ref{thm:one3}
forces a strict drop in $v_{3}$ at $i_{0}+1$, and $v_{3}$ never increases along an
orbit (\texttt{v3\_orbit\_tail}); returning to $d$ at step $L$ is then impossible.
So $3\nmid d$.  With the prime $3$ absent, Lemma~\ref{thm:gcd3} forces $g=1$ at
every step, so the denominator is constant along the whole cycle.
\end{proof}

\subsection{Descent on pure powers of three}

Let \texttt{oddCount x n} count odd-numerator steps in the first $n$ steps of
the orbit of $x$.  The main theorem is:

\begin{theorem}[Theorem 1]
\label{thm:descent}
\begin{lstlisting}
theorem theorem_1_pure_three_power_descent (a k : ℕ) (ha : 0 < a)
    (hcop : Nat.Coprime a (3 ^ k)) :
    ∃ t m, step^[t] (a, 3 ^ k) = (m, 1) ∧ oddCount (a, 3 ^ k) t ≤ k
\end{lstlisting}
\end{theorem}

\begin{proof}[Proof outline]
Induction on $k$.  Factor $a=2^{s}b$ with $b$ odd (\texttt{even\_steps\_to\_odd});
the first $s$ steps are halvings, so they leave the denominator $3^{k+1}$
unchanged and contribute no odd steps.  The next step is odd; by
Lemma~\ref{thm:one3} its gcd with $3^{k+1}$ is exactly $3$, so it sends
$(b,3^{k+1})$ to $(b+3^{k},3^{k})$.  The induction hypothesis applied to the new
numerator, whose coprimality with $3^{k}$ is inherited, supplies the remaining
$t'$ steps and the count bound.  Additivity of \texttt{oddCount}
(\texttt{oddCount\_add}) closes the arithmetic.
\end{proof}

\subsection{Exhaustive verification for \texorpdfstring{$1\le k\le 6$}{1 <= k <= 6}}

The Boolean sieve and its reflected propositional form are:

\begin{lstlisting}
def reachesOneWithin (k a B : ℕ) : Prop :=
  ∃ t ≤ B, step^[t] (a, 3 ^ k) = (1, 1)

def reachesOneBool (k a B : ℕ) : Bool :=
  (List.range (B + 1)).any (fun t => step^[t] (a, 3 ^ k) == (1, 1))

def verifyPureThreePowerSmallK : Bool :=
  (List.range 6).all fun ki =>
    (List.range (3 ^ (ki + 1))).all fun ai =>
      if Nat.Coprime ai (3 ^ (ki + 1)) then reachesOneBool (ki + 1) ai 100 else true

set_option maxRecDepth 100000 in
set_option maxHeartbeats 800000 in
lemma verification_pure_three_power_k_le_six : verifyPureThreePowerSmallK = true := by
  decide

lemma verification_all_reduced_small_k :
    ∀ k : Fin 6, ∀ a : Fin (3 ^ (k.val + 1)),
      Nat.Coprime a.val (3 ^ (k.val + 1)) →
      reachesOneWithin (k.val + 1) a.val 100
\end{lstlisting}

\noindent The search bound $B=100$ is deliberately conservative: direct
evaluation of the same reduced dynamics over all reduced numerators
$0<a<3^{k}$ for $1\le k\le 6$ shows that the true maximum number of reduced
steps is \textbf{35}, attained at $k=6$, $a=520$.  Thus the kernel check certifies
a bound with a comfortable margin, and the observed extremal trajectory fits
well inside it.  The propositional form quantifies over all $\sum_{k=1}^{6}\phi(3^{k})$
reduced states in the range, i.e.\ $2+6+18+54+162+486=728$ states.

\section{Mechanized 2-adic fuel and alternating-bit bounds}
\label{sec:glide}

This section summarizes two modules.  The first,
\texttt{CollatzAdicBase/FuelLaw.lean}, proves the fuel law; the second,
\texttt{CollatzAdicBase/AlternatingGlide.lean}, proves the glide quantization.
Both glide predicates were given in Section~\ref{sec:arch}.

\subsection{The initial fuel law}

The algebraic core is the exact evolution of the shortcut orbit:

\begin{lstlisting}
theorem iterate_shortcut_formula (k q j : Nat) (hk : 1 ≤ k) (hq : 1 ≤ q) (hj : j ≤ k) :
    Nat.iterate shortcut j (2 ^ k * q - 1) = 3 ^ j * 2 ^ (k - j) * q - 1

theorem collatz_iterate_even (k q j : Nat) (hk : 1 ≤ k) (hq : 1 ≤ q) (hj : j ≤ k) :
    Nat.iterate collatz (2 * j) (2 ^ k * q - 1) = 3 ^ j * 2 ^ (k - j) * q - 1

theorem collatz_iterate_two_mul (k q : Nat) (hk : 1 ≤ k) (hq : 1 ≤ q) :
    Nat.iterate collatz (2 * k) (2 ^ k * q - 1) = 3 ^ k * q - 1
\end{lstlisting}

\noindent The strict ascent for all intermediate times, and the resulting glide
bound, are:

\begin{proposition}[Proposition 4]
\label{thm:fuel}
\begin{lstlisting}
theorem collatz_iterate_gt_initial (k q t : Nat) (hk : 1 ≤ k) (hq : 1 ≤ q)
    (ht1 : 1 ≤ t) (ht2 : t ≤ 2 * k) :
    Nat.iterate collatz t (2 ^ k * q - 1) > 2 ^ k * q - 1

theorem glide_at_least_two_mul (k q : Nat) (hk : 1 ≤ k) (hq : 1 ≤ q) :
    GlideAtLeast (2 ^ k * q - 1) (2 * k)
\end{lstlisting}
\end{proposition}

\begin{proof}[Proof outline]
Two ordinary steps from an odd $m$ equal one shortcut step
(\texttt{collatz\_twice\_odd}), so after $2j$ steps the value is the $j$-th
shortcut value above.  For $1\le j\le k$ one has $2^{j}<3^{j}$ and
$2^{k}q\le 3^{j}2^{k-j}q$, which yields strict inequality after subtracting $1$;
the odd intermediate times are handled by one additional odd step.  The endpoint
formula at $j=k$ gives $\texttt{collatz}^{2k}(2^{k}q-1)=3^{k}q-1$.
\end{proof}

Since the orbit exceeds its starting value for every $t$ with $1\le t\le 2k$,
the first glide time satisfies $G(2^{k}q-1)\ge 2k+1$.  This is the formal
counterpart of the $2$-adic fuel law of~\cite{paper2}.

\subsection{Alternating-bit glide quantization}

For $k\ge 1$ define the alternating-bit family
\[
  n_{\mathrm{alt}}(k)=\frac{2^{2k+3}+1}{3},
\]
whose binary expansion is $10\,10\cdots 10\,11$ with $k$ copies of $10$.  The
definition is well-defined in $\mathbb{N}$ because $2^{2k+3}\equiv 2\pmod 3$:

\begin{lstlisting}
def n_alt (k : Nat) : Nat := (2 ^ (2 * k + 3) + 1) / 3

theorem three_dvd_alt (k : Nat) : 3 ∣ 2 ^ (2 * k + 3) + 1

theorem three_mul_n_alt (k : Nat) : 3 * n_alt k = 2 ^ (2 * k + 3) + 1
\end{lstlisting}

\noindent Writing $r=2^{2k-2}$, an eight-step table is proved step by step:

\begin{lstlisting}
theorem n_alt_iterate_two (k : Nat) (hk : 1 ≤ k) :
    Nat.iterate collatz 2 (n_alt k) = 2 ^ (2 * k + 2) + 1

theorem n_alt_iterate_five (k : Nat) (hk : 1 ≤ k) :
    Nat.iterate collatz 5 (n_alt k) = 3 * 2 ^ (2 * k) + 1

theorem n_alt_iterate_eight (k : Nat) (hk : 1 ≤ k) :
    Nat.iterate collatz 8 (n_alt k) = 9 * 2 ^ (2 * k - 2) + 1
\end{lstlisting}

\begin{theorem}[Theorem 1 of \texttt{AlternatingGlide.lean}]
\label{thm:alt}
\begin{lstlisting}
theorem n_alt_glide_eq (k : Nat) (hk : 1 ≤ k) : GlideEq (n_alt k) 8

theorem alternating_bit_glide (k : Nat) (hk : 1 ≤ k) :
    Nat.iterate collatz 2 (n_alt k) = 2 ^ (2 * k + 2) + 1 ∧
    Nat.iterate collatz 5 (n_alt k) = 3 * 2 ^ (2 * k) + 1 ∧
    Nat.iterate collatz 8 (n_alt k) = 9 * 2 ^ (2 * k - 2) + 1 ∧
    (∀ t, 1 ≤ t → t ≤ 7 → Nat.iterate collatz t (n_alt k) > n_alt k) ∧
    Nat.iterate collatz 8 (n_alt k) < n_alt k
\end{lstlisting}
\end{theorem}

\begin{proof}[Proof outline]
From $3n_{\mathrm{alt}}(k)=32r+1$ one obtains the closed forms of
$T^{t}(n_{\mathrm{alt}}(k))$ for $t=1,\dots,8$ by eight explicit lemmas, each
discharging a parity case with \texttt{omega}.  For $1\le t\le 7$ the value
exceeds $9r+3$ while $3n_{\mathrm{alt}}(k)=32r+1$, so $T^{t}>n_{\mathrm{alt}}$;
at $t=8$ the value is $9r+1<10r+\tfrac{1}{3}=n_{\mathrm{alt}}$.  This gives
\texttt{GlideEq} and hence $G(n_{\mathrm{alt}}(k))=8$.
\end{proof}

\section{Verification-driven corrigendum: \texorpdfstring{$n_{0}=3$}{n0 = 3}}
\label{sec:corrigendum}

Lemma 3 of~\cite{paper2} concerns the exceptional case $k=0$, where
$n_{\mathrm{alt}}(0)=3$.  The printed statement asserts
$T^{6}(n_{0})=1<n_{0}$ and $T^{8}(n_{0})=4>n_{0}$.  The second equality is
correct; the first is off by one.  Under the map actually formalized in
\texttt{Basic.lean}, the orbit of $3$ is kernel-checked by
\texttt{lemma\_three\_orbit} and displayed in full below.

\begin{center}
\begin{tabular}{@{}cccccccccc@{}}
\toprule
$t$ & 0 & 1 & 2 & 3 & 4 & 5 & 6 & 7 & 8 \\
\midrule
$T^{t}(3)$ & 3 & 10 & 5 & 16 & 8 & 4 & \textbf{2} & \textbf{1} & \textbf{4} \\
\bottomrule
\end{tabular}
\end{center}

\noindent The corresponding verified statement is

\begin{lstlisting}
theorem lemma_three_orbit :
    Nat.iterate collatz 6 3 = 2 ∧
    Nat.iterate collatz 7 3 = 1 ∧
    Nat.iterate collatz 8 3 = 4 := by
  decide

theorem lemma_three_k_zero :
    n_alt 0 = 3 ∧
    Nat.iterate collatz 7 3 = 1 ∧ Nat.iterate collatz 7 3 < 3 ∧
    Nat.iterate collatz 8 3 = 4 ∧ Nat.iterate collatz 8 3 > 3
\end{lstlisting}

so that
\[
  T^{6}(3)=2,\qquad T^{7}(3)=1<3,\qquad T^{8}(3)=4>3 .
\]
The total stopping time of $3$ is therefore $7$ applications of $T$, not $6$.
The qualitative claim of the printed lemma survives: the orbit has already
dropped below its start by step $6$ ($T^{6}(3)=2<3$); only the index attached to
the value $1$ was misstated.  Theorem~\ref{thm:alt} (valid for $k\ge 1$) and
Proposition~\ref{thm:fuel} are unaffected: for $k\ge 1$ one has
$n_{\mathrm{alt}}(k)\ge 11$ and $G(n_{\mathrm{alt}}(k))=8$, and the $k=0$ case is
the isolated exception recorded by Lemma 3.  A corrigendum is archived with the
source release as \texttt{CORRIGENDUM.md}; no revised version of~\cite{paper2}
has been deposited on Zenodo at the time of writing, so that the corrected PDF
and this paper can be cross-linked directly.

\section{Axiom audit and reproducibility}
\label{sec:audit}

\subsection{The axiom audit}

\texttt{CollatzAdicBase/Axioms.lean} issues \texttt{\#print axioms} for 29
representative theorems across the five verified modules.  Running

\begin{lstlisting}
lake env lean CollatzAdicBase/Axioms.lean > AXIOM_AUDIT.txt 2>&1
\end{lstlisting}

\noindent produces the following verbatim table, also archived as
\texttt{AXIOM\_AUDIT.txt}:

\begin{lstlisting}
'CollatzAdicBase.FuelLaw.iterate_shortcut_formula' depends on axioms: [propext, Quot.sound]
'CollatzAdicBase.FuelLaw.iterate_shortcut_odd' depends on axioms: [propext, Quot.sound]
'CollatzAdicBase.FuelLaw.iterate_shortcut_strict_mono' depends on axioms: [propext, Quot.sound]
'CollatzAdicBase.FuelLaw.collatz_iterate_gt_initial' depends on axioms: [propext, Quot.sound]
'CollatzAdicBase.FuelLaw.collatz_iterate_two_mul' depends on axioms: [propext, Quot.sound]
'CollatzAdicBase.FuelLaw.glide_at_least_two_mul' depends on axioms: [propext, Quot.sound]
'CollatzAdicBase.lemma_2_1_stability' depends on axioms: [propext, Classical.choice, Quot.sound]
'CollatzAdicBase.lemma_2_2_one_three_per_odd_step' depends on axioms: [propext, Classical.choice, Quot.sound]
'CollatzAdicBase.lemma_2_3_forced_halving' depends on axioms: [propext, Classical.choice, Quot.sound]
'CollatzAdicBase.lemma_2_4_and_cor_2_5_three_free_cycles' depends on axioms: [propext, Classical.choice, Quot.sound]
'CollatzAdicBase.theorem_1_pure_three_power_descent' depends on axioms: [propext, Classical.choice, Quot.sound]
'CollatzAdicBase.verification_pure_three_power_k_le_six' depends on axioms: [propext, Classical.choice, Quot.sound]
'CollatzAdicBase.verification_all_reduced_small_k' depends on axioms: [propext, Classical.choice, Quot.sound]
'CollatzAdicBase.Refueling.two_pow_dvd_iff_le_v2' depends on axioms: [propext, Quot.sound]
'CollatzAdicBase.Refueling.refuel_landing' depends on axioms: [propext, Quot.sound]
'CollatzAdicBase.Refueling.refuel_exhausts_fuel' depends on axioms: [propext, Quot.sound]
'CollatzAdicBase.Refueling.refuel_congruence' depends on axioms: [propext, Quot.sound]
'CollatzAdicBase.Refueling.refuel_mod_invariant' depends on axioms: [propext, Quot.sound]
'CollatzAdicBase.AlternatingGlide.three_dvd_alt' depends on axioms: [propext, Quot.sound]
'CollatzAdicBase.AlternatingGlide.n_alt_glide_eq' depends on axioms: [propext, Quot.sound]
'CollatzAdicBase.AlternatingGlide.alternating_bit_glide' depends on axioms: [propext, Quot.sound]
'CollatzAdicBase.AlternatingGlide.lemma_three_k_zero' depends on axioms: [propext, Classical.choice, Quot.sound]
'CollatzAdicBase.LatentSaturation.gcd_lift' depends on axioms: [propext, Quot.sound]
'CollatzAdicBase.LatentSaturation.ratStep_lift' depends on axioms: [propext, Quot.sound]
'CollatzAdicBase.LatentSaturation.lift_congruence' depends on axioms: [propext]
'CollatzAdicBase.LatentSaturation.cycle_dvd_pow_sub' depends on axioms: [propext, Quot.sound]
'CollatzAdicBase.LatentSaturation.den_at_lifted_level' depends on axioms: [propext, Quot.sound]
'CollatzAdicBase.LatentSaturation.three_free_preserved' depends on axioms: [propext, Quot.sound]
'CollatzAdicBase.LatentSaturation.ratStep_valid' depends on axioms: [propext, Quot.sound]
\end{lstlisting}

\noindent Every entry is a subset of $\{\texttt{propext},\texttt{Classical.choice},
\texttt{Quot.sound}\}$.  In particular \texttt{sorryAx} does not appear, and no
axiom declared by the project appears: the development is conservative over
Lean's kernel.  The presence of \texttt{Classical.choice} in the rational-descent
theorems is routine, arising from classical reasoning about coprimality and the
use of \texttt{Nat.find}; the $2$-adic modules avoid even that dependency.

\subsection{Build instructions and release information}

The complete verification is reproduced by the following commands from the
repository root:

\begin{lstlisting}
# Toolchain (pinned in ./lean-toolchain): leanprover/lean4:v4.36.0
lake build
lake env lean CollatzAdicBase/Axioms.lean > AXIOM_AUDIT.txt 2>&1
\end{lstlisting}

\noindent Mathlib is pinned to \texttt{v4.36.0} in \texttt{lakefile.toml} and
\texttt{lake-manifest.json}.  A full clean rebuild of the library reports
\texttt{Build completed successfully (18266 jobs)} with zero errors and zero
warnings.  The unified state described here is the commit carrying the
annotated tag \texttt{v1.0.0}.

\section{Conclusion}

We have given a machine-checked account of the local valuation dynamics used in
the first two preprints of this series.  The deliverables are (i) a verified
$3$-adic descent framework in which cycles are proved $3$-free and pure
$3$-power denominators reduce to integers, with a kernel-checked exhaustive
check for $1\le k\le 6$; (ii) verified $2$-adic glide bounds, including the
exact quantization $G(n_{\mathrm{alt}}(k))=8$ for all $k\ge 1$; and (iii) a
corrected Lemma 3 for the exceptional input $n_{0}=3$, including the full
$T$-orbit table.  The axiom audit certifies that nothing beyond Lean's core
axioms is assumed.  The accompanying Lean repository, released as \texttt{v1.0.0},
contains all sources, the corrigendum, and the archived audit.

\paragraph{Data availability.}
The two prior preprints are archived on Zenodo with DOIs
\href{https://doi.org/10.5281/zenodo.24375499}{10.5281/zenodo.24375499}
(Paper 1) and
\href{https://doi.org/10.5281/zenodo.24397999}{10.5281/zenodo.24397999}
(Paper 2).  The verified Lean sources and the artifact
\texttt{AXIOM\_AUDIT.txt} accompany the present paper in the tagged release.

\begin{thebibliography}{9}

\bibitem{paper1} J.~W. Simon.
\newblock Periodic cycles of the 2-adic rational Collatz map: 3-adic exclusivity
  and descent on pure 3-power denominators.
\newblock Zenodo preprint, 2026.
\newblock \href{https://doi.org/10.5281/zenodo.24375499}{doi:10.5281/zenodo.24375499}.

\bibitem{paper2} J.~W. Simon.
\newblock Trailing Ones and High-Altitude Backbones: Empirical Trajectory
  Dynamics and Algebraic Sub-Family Bounds in the Integer Collatz Map up to
  $10^{7}$.
\newblock Zenodo preprint, 2026.
\newblock \href{https://doi.org/10.5281/zenodo.24397999}{doi:10.5281/zenodo.24397999}.

\bibitem{terras} R.~Terras.
\newblock A stopping time problem on the positive integers.
\newblock \emph{Acta Arithmetica}, 30:241--252, 1976.

\bibitem{lagarias} J.~C.~Lagarias, editor.
\newblock \emph{The $3x+1$ Problem and Its Generalizations}.
\newblock American Mathematical Society, 2010.

\bibitem{tao} T.~Tao.
\newblock Almost all orbits of the Collatz map attain almost bounded values.
\newblock \emph{Forum of Mathematics, Pi}, 10:e12, 2022.

\bibitem{mathlib} The mathlib community.
\newblock The Lean mathematical library.
\newblock In \emph{Proceedings of the 9th ACM SIGPLAN International Conference on
  Certified Programs and Proofs (CPP 2020)}, pages 367--381, 2020.

\bibitem{lean4} L.~de Moura and S.~Ullrich.
\newblock The Lean 4 theorem prover and programming language.
\newblock In \emph{Automated Deduction -- CADE 28}, LNCS 12699, pages 625--635,
  Springer, 2021.

\end{thebibliography}

\end{document}
